% Generated draft: requires mathematical and source review.
\documentclass[11pt]{article}
\usepackage[margin=25mm]{geometry}
\usepackage{fontspec}
\setmainfont{Latin Modern Roman}
\usepackage{amsmath,amssymb,mathtools,unicode-math}
\setmathfont{Latin Modern Math}
\usepackage{enumitem,needspace,xurl,hyperref}
\setlength{\parindent}{0pt}
\setlength{\parskip}{5pt}
\setlength{\emergencystretch}{3em}
\newcommand{\N}{\mathbb{N}}
\newcommand{\Z}{\mathbb{Z}}
\newcommand{\Q}{\mathbb{Q}}
\newcommand{\R}{\mathbb{R}}
\newcommand{\C}{\mathbb{C}}
\newcommand{\F}{\mathbb{F}}
\newcommand{\A}{\mathbb{A}}
\newcommand{\PP}{\mathbb{P}}
\newcommand{\E}{\mathbb{E}}
\newcommand{\eps}{\varepsilon}
\newcommand{\dd}{\,\mathrm{d}}
\begin{document}
\section*{UnsolvedMath ID: 800620}
\textbf{Source problem code:} catalogue entry
\subsection{Definitions and mathematical statement}
Under the conventional interpretation of the source notation, write $[\N]^{\aleph_0}$ for the family of infinite subsets of $\N$, and write $A\subseteq^*B$ when $A\setminus B$ is finite. A family $\mathcal F\subseteq[\N]^{\aleph_0}$ has the strong finite intersection property if every nonempty finite subfamily has infinite intersection. An infinite pseudointersection of $\mathcal F$ is a set $A\in[\N]^{\aleph_0}$ such that $A\subseteq^*F$ for every $F\in\mathcal F$. The pseudointersection number $\mathfrak p$ is the least cardinality of a family with the strong finite intersection property and no infinite pseudointersection. The tower number $\mathfrak t$ is the least cardinal $\kappa$ for which there is a family $(T_\alpha)_{\alpha<\kappa}$ of infinite subsets of $\N$ such that $T_\beta\subseteq^*T_\alpha$ whenever $\alpha<\beta<\kappa$, and the family has no infinite pseudointersection. These are conventional definitions; the supplied excerpt does not itself define these symbols.

The source also uses cardinal invariants $\mathfrak p_\kappa$, but the supplied record gives no definition of them. It reports the identities
\[
\mathfrak p=\mathfrak p_2=\mathfrak p_3=\cdots=\mathfrak p_{\aleph_0},
\qquad
\mathfrak p_{\mathfrak t}=\mathfrak t,
\]
and attributes $\mathfrak p=\mathfrak p_{\aleph_1}$ to the reference label [Pyt]. It displays the hierarchy
\[
\mathfrak p=\mathfrak p_{\aleph_1}\leq\mathfrak p_{\aleph_2}\leq\cdots\leq\mathfrak p_{\mathfrak t}=\mathfrak t.
\]
These are supplied background assertions, not independently verified conclusions of this draft.

The question in Problem 49 is whether
\[
\mathfrak p_{\mathfrak p}=\mathfrak t.
\]
Here $\mathfrak p_{\mathfrak p}$ denotes the source's indexed invariant evaluated at the cardinal $\mathfrak p$. A fully specified mathematical statement requires the missing definition of $\mathfrak p_\kappa$, including its admissible indices and defining quantifiers. The reported identities and inequalities do not supply that definition. The underlying axiomatic assumptions are also not stated in the supplied excerpt.
\subsection{Short English statement}
Tsaban's Problem 49 asks whether the source's cardinal invariant $\mathfrak p_\kappa$, evaluated at $\kappa=\mathfrak p$, equals the tower number $\mathfrak t$. The supplied record labels the problem open, but omits the definition of $\mathfrak p_\kappa$, so the formulation remains incomplete and requires human mathematical review.
\subsection{Sources}
\begin{itemize}
\item Local UnsolvedMath registry, permanent ID 800620, exact source problem code \texttt{catalogue entry}, title ``Selection Principles in Mathematics --- Problem 49.'' The supplied record attributes the problem to Tsaban (2006) and labels its status ``open''; this label is supplied metadata, not an independently verified current assessment.
\item Tsaban, \emph{Selection Principles and Special Sets of Reals: Open Problems} (2006), Problem 49. Supplied source URL: \url{https://arxiv.org/abs/math/0609601}. The provenance records access on 2026-10-16 and extraction by \texttt{arxiv\_latex}; these are supplied provenance claims, not a claim of access or verification for this draft.
\item The supplied excerpt cites [Pyt] for $\mathfrak p=\mathfrak p_{\aleph_1}$, but supplies no bibliographic details identifying that reference.
\end{itemize}
\end{document}
